% Save this file as UTF-8. Compile with XeLaTeX.
\documentclass[10pt,a4paper]{article}
\usepackage[margin=1.65cm]{geometry}
\usepackage{fontspec}
\setmainfont{TeX Gyre Heros}
\usepackage{enumitem}
\usepackage{titlesec}
\usepackage[hidelinks]{hyperref}
\pagestyle{empty}
\setlength{\parindent}{0pt}
\setlength{\parskip}{3pt}
\setlist[itemize]{leftmargin=1.25em,itemsep=2pt,topsep=3pt,parsep=0pt}
\titleformat{\section}{\large\bfseries}{}{0pt}{}[\titlerule]
\titlespacing*{\section}{0pt}{10pt}{5pt}

\begin{document}
\begin{center}
  {\small CURRICULUM VITAE}\\[4pt]
  {\LARGE\bfseries QI WANG}\\[4pt]
  Master's Student | Shanghai University\\[3pt]
  +86 199 0322 5165\quad |\quad
  \href{mailto:1985277475@qq.com}{1985277475@qq.com}\quad |\quad Shanghai, China
\end{center}

\section*{EDUCATION}
\textbf{Shanghai University}\hfill Sep 2026 -- Present\\
M.Eng. Electronic Information\\
Research interests: evolutionary learning and medical image segmentation

\textbf{Civil Aviation University of China}\hfill Sep 2021 -- Jul 2025\\
B.Eng. Electrical Engineering and Automation | GPA: 3.71 | Top 10\%\\
Selected coursework: Signals and Systems (96), Control Theory and Engineering (90),
C Programming (95), Power Electronics (92), Electrical Drive Control Systems (92)

\section*{RESEARCH \& ENGINEERING PROJECTS}
\textbf{Human Learning Optimisation for Medical Image Segmentation}
\begin{itemize}
  \item Investigated HLO-based neural architecture and parameter optimisation for medical image segmentation, adopting TransUNet as the baseline model
  \item Reviewed and refactored source code of HLO variants to support algorithm integration and experimental validation on medical datasets
\end{itemize}

\textbf{Aircraft Longitudinal Flight Control | Bachelor's Thesis}
\begin{itemize}
  \item Developed and linearised a longitudinal dynamics model with hierarchical pitch, altitude and velocity control
  \item Implemented online PID gain tuning using a 49-rule fuzzy inference system driven by tracking error and its derivative
  \item MATLAB simulations demonstrated faster transient response and superior steady-state accuracy compared with fixed-gain PID under step inputs and disturbances
\end{itemize}

\textbf{Aircraft Attitude Control during Three-Point Jacking}
\begin{itemize}
  \item Developed an STM32/MPU6050 acquisition system with I\textsuperscript{2}C communication, complementary filtering and quaternion-based attitude estimation
  \item Modelled three-point jacking dynamics with PID feedback and SVPWM vector control of a permanent-magnet synchronous motor
  \item Achieved $\pm 0.3^{\circ}$ pitch and roll accuracy in simulation, validated via static calibration and step-response tests
\end{itemize}

\section*{INDUSTRIAL EXPERIENCE}
\textbf{Shanghai Baoying Food \& Beverage Co., Ltd.}\\
\textit{Quality Control Intern}\hfill Jun--Aug 2023; Jan--Feb 2024
\begin{itemize}
  \item Monitored production and product quality; documented non-conformities for statistical quality analysis
  \item Collaborated with production teams on formulation verification, sampling, water quality testing and abnormal quality event handling
\end{itemize}

\section*{SKILLS \& AWARDS}
\textbf{Programming \& Tools:} Python, MATLAB, C, STM32, I\textsuperscript{2}C, Visio, MS Office\\
\textbf{Domain Expertise:} Evolutionary learning optimisation, medical image segmentation, PID \& fuzzy control, attitude estimation\\
\textbf{Languages:} English (CET-4)\\
\textbf{Awards:} People's Scholarship (Second Class)
\end{document}
